La planificación con consignas sin pistas y criterios de éxito observables hizo comparables las tareas T1 a T3. La evaluación heurística y el recorrido cognitivo anticiparon problemas de recuperación y visibilidad que las correcciones v2 atienden, y la auditoría de accesibilidad mostró un flujo operable por teclado con un contraste mejorable en el aviso verde. La eficacia, la eficiencia y la satisfacción se describen para tres participantes, sin generalizar.
